% GENERATED FILE. Edit canonical lessons, then run npm run generate:books.
% canonical-source-hash: fnv1a64:c900d73f69c8a4b7
% canonical-lessons: SA-C245-nalika, SA-C245-takram, SA-C245-phalarasah, SA-C245-lavangam, SA-C245-ela

\chapter{Tube, Buttermilk, Fruit juice, Clove, Cardamom}
\label{ch:sa-tube-buttermilk-fruit-juice-clove-cardamom}
\hlchaptermodality{245}

\begin{chapteropening}
\textbf{By the end of this chapter:} \emph{I can say a tube, buttermilk, fruit juice, a clove and cardamom.}

Complete the last lesson of chapter 245: I can say a tube, buttermilk, fruit juice, a clove and cardamom.
\end{chapteropening}

\section[\texorpdfstring{\sk{नालिका} (nālikā)}{नालिका (nālikā)}]{\sk{नालिका} (nālikā) — a tube, a pipe}

\label{lesson:SA-C245-nalika}



\begin{quote}
Before the new one: say the Sanskrit for other, then the Sanskrit for alone.
\end{quote}

\subsection*{You'll want to know: \sk{नालिका}}
\textbf{\sk{नालिका}} — \emph{nālikā} — \textquotedblleft{}a tube, a pipe\textquotedblright{}.

\begin{grammarlens}[title={What you've built}]
One more word for this chapter.
\end{grammarlens}

\subsection*{Your turn}
\begin{itemize}
\raggedright
  \item \emph{Say it:} \emph{nālikā}
  \item \emph{Say it:} \emph{nālikā}, once more
  \item \emph{Say it:} say \emph{nālikā}
  \item \emph{Recall:} say the Sanskrit for other, then the Sanskrit for alone, then say \emph{nālikā} again
\end{itemize}

\begin{tcolorbox}[breakable,colback=teal!4,colframe=teal!35!black,title={Before you move on}]
What does \sk{नालिका} mean? (\textquotedblleft{}A tube, a pipe\textquotedblright{}.) Say it once more.
\end{tcolorbox}
\section[\texorpdfstring{\sk{तक्रम्} (takram)}{तक्रम् (takram)}]{\sk{तक्रम्} (takram) — buttermilk}

\label{lesson:SA-C245-takram}



\begin{quote}
Before the new one: say the Sanskrit for alone, then the Sanskrit for a tube.
\end{quote}

\subsection*{You'll want to know: \sk{तक्रम्}}
\textbf{\sk{तक्रम्}} — \emph{takram} — \textquotedblleft{}buttermilk\textquotedblright{}.

\begin{grammarlens}[title={What you've built}]
One more word for this chapter.
\end{grammarlens}

\subsection*{Your turn}
\begin{itemize}
\raggedright
  \item \emph{Say it:} \emph{takram}
  \item \emph{Say it:} \emph{takram}, once more
  \item \emph{Say it:} say \emph{takram}
  \item \emph{Recall:} say the Sanskrit for alone, then the Sanskrit for a tube, then say \emph{takram} again
\end{itemize}

\begin{tcolorbox}[breakable,colback=teal!4,colframe=teal!35!black,title={Before you move on}]
What does \sk{तक्रम्} mean? (\textquotedblleft{}Buttermilk\textquotedblright{}.) Say it once more.
\end{tcolorbox}
\section[\texorpdfstring{\sk{फलरसः} (phalarasaḥ)}{फलरसः (phalarasaḥ)}]{\sk{फलरसः} (phalarasaḥ) — fruit juice}

\label{lesson:SA-C245-phalarasah}



\begin{quote}
Before the new one: say the Sanskrit for a tube, then the Sanskrit for buttermilk.
\end{quote}

\subsection*{You'll want to know: \sk{फलरसः}}
\textbf{\sk{फलरसः}} — \emph{phalarasaḥ} — \textquotedblleft{}fruit juice\textquotedblright{}.

\begin{grammarlens}[title={What you've built}]
One more word for this chapter.
\end{grammarlens}

\subsection*{Your turn}
\begin{itemize}
\raggedright
  \item \emph{Say it:} \emph{phalarasaḥ}
  \item \emph{Say it:} \emph{phalarasaḥ}, once more
  \item \emph{Say it:} say \emph{phalarasaḥ}
  \item \emph{Recall:} say the Sanskrit for a tube, then the Sanskrit for buttermilk, then say \emph{phalarasaḥ} again
\end{itemize}

\begin{tcolorbox}[breakable,colback=teal!4,colframe=teal!35!black,title={Before you move on}]
What does \sk{फलरसः} mean? (\textquotedblleft{}Fruit juice\textquotedblright{}.) Say it once more.
\end{tcolorbox}
\section[\texorpdfstring{\sk{लवंगम्} (lavaṅgam)}{लवंगम् (lavaṅgam)}]{\sk{लवंगम्} (lavaṅgam) — a clove}

\label{lesson:SA-C245-lavangam}



\begin{quote}
Before the new one: say the Sanskrit for buttermilk, then the Sanskrit for fruit juice.
\end{quote}

\subsection*{You'll want to know: \sk{लवंगम्}}
\textbf{\sk{लवंगम्}} — \emph{lavaṅgam} — \textquotedblleft{}a clove\textquotedblright{}.

\begin{grammarlens}[title={What you've built}]
One more word for this chapter.
\end{grammarlens}

\subsection*{Your turn}
\begin{itemize}
\raggedright
  \item \emph{Say it:} \emph{lavaṅgam}
  \item \emph{Say it:} \emph{lavaṅgam}, once more
  \item \emph{Say it:} say \emph{lavaṅgam}
  \item \emph{Recall:} say the Sanskrit for buttermilk, then the Sanskrit for fruit juice, then say \emph{lavaṅgam} again
\end{itemize}

\begin{tcolorbox}[breakable,colback=teal!4,colframe=teal!35!black,title={Before you move on}]
What does \sk{लवंगम्} mean? (\textquotedblleft{}A clove\textquotedblright{}.) Say it once more.
\end{tcolorbox}
\section[\texorpdfstring{\sk{एला} (elā)}{एला (elā)}]{\sk{एला} (elā) — cardamom}

\label{lesson:SA-C245-ela}



\begin{quote}
Before the new one: say the Sanskrit for fruit juice, then the Sanskrit for a clove.
\end{quote}

\subsection*{You'll want to know: \sk{एला}}
\textbf{\sk{एला}} — \emph{elā} — \textquotedblleft{}cardamom\textquotedblright{}.

\begin{grammarlens}[title={What you've built}]
One more word for this chapter.
\end{grammarlens}

\subsection*{Your turn}
\begin{itemize}
\raggedright
  \item \emph{Say it:} \emph{elā}
  \item \emph{Say it:} \emph{elā}, once more
  \item \emph{Say it:} say \emph{elā}
  \item \emph{Recall:} say the Sanskrit for fruit juice, then the Sanskrit for a clove, then say \emph{elā} again
\end{itemize}

\begin{tcolorbox}[breakable,colback=teal!4,colframe=teal!35!black,title={Before you move on}]
What does \sk{एला} mean? (\textquotedblleft{}Cardamom\textquotedblright{}.) Say it once more.
\end{tcolorbox}
